Lors d'une randonnée en montagne, un alpiniste, de masse $m=\qty{80}{\kg}$, part
de Chamonix (altitude $\qty{1050}{\m}$). Il passe au sommet du Mont Blanc
(altitude $\qty{4807}{\m}$) et redescend à Megève (altitude $\qty{1113}{\m}$).

Calculer le travail du poids de l'alpiniste pendant cette randonnée. Commenter
le résultat.

\begin{donnees}
	$g=\qty{10}{\N\per\kg}$.
\end{donnees}
